%% ──────────────────────────────────────────────
%%  فصل ۱۵ – بالاتر از سیاهی رنگی نیست:
%%  نجات‌طلبی فداکارانه و مسئولیت نسلی
%%  فایل: chapters/ch15-sacrificial-liberation.tex
%% ──────────────────────────────────────────────
\chapter{بالاتر از سیاهی رنگی نیست: نجات‌طلبی فداکارانه و مسئولیت نسلی}
\label{ch:sacrificial-liberation}

%% ── چکیده‌ی فصل ──
\begin{tcolorbox}[mybox, title={چکیده‌ی فصل}]
فصل‌های پیشین این کتاب استدلال کردند که نجات‌طلبی — امیدبستن به نیروی خارجی برای رهایی — در اکثر موارد تاریخی شکست خورده و نتایجی از بی‌ثباتی تا فاجعه به‌بار آورده است. اما آیا این تحلیل، تمام داستان است؟ این فصل عمداً نقش \concept{مدافع شیطان} (\lat{Devil's Advocate}) را ایفا می‌کند و قوی‌ترین استدلال‌های ممکن را \textbf{به نفع} نجات‌طلبی — آنچه \concept{نجات‌طلبی فداکارانه} می‌نامیم — ارائه می‌دهد. پرسش محوری این است: \textbf{وقتی انسداد سیاسی مطلق است و هیچ راه داخلی‌ای وجود ندارد، آیا پذیرش هزینه‌ی مداخله‌ی خارجی — حتی به قیمت یک یا دو نسل — برای رهایی نسل‌های آینده موجه نیست؟} اگر «بالاتر از سیاهی رنگی نیست»، آیا هر تغییری — حتی پرهزینه — بهتر از تداوم وضع موجود نیست؟ این فصل ابتدا چارچوب‌های مختلف پاسخ‌گویی به این پرسش را معرفی می‌کند، سپس شواهد تاریخی و نظری هر رویکرد را بررسی و در نهایت ارزیابی می‌نماید.
\end{tcolorbox}

%% ================================================================
\section{درآمد: پرسشی که نباید نادیده گرفت}
\label{sec:ch15-intro}
%% ================================================================

\needspace{6\baselineskip}
نویسنده‌ی این کتاب در فصل‌های ۱ تا ۱۴ استدلال کرد که نجات‌طلبی معمولاً آرزواندیشی‌ای خطرناک است. اما صداقت فکری ایجاب می‌کند که قوی‌ترین استدلال مخالف نیز شنیده شود — استدلالی که از دل تاریک‌ترین تجربه‌های بشری برمی‌آید:

\begin{tcolorbox}[quotebox, title={صدای دیگر}]
«شما از نروژ یا سوئیس درباره‌ی عاملیت جمعی و مقاومت مدنی سخن می‌گویید. اما من از آشویتس سخن می‌گویم. از حلبچه. از سربرنیتسا. از رواندا. در آشویتس هیچ مقاومت مدنی‌ای کارساز نبود. هیچ نهاد مدنی مستقلی نمی‌توانست کوره‌های آدم‌سوزی را متوقف کند. تنها ارتش سرخ و متفقین توانستند آشویتس را آزاد کنند. آیا یهودیانی که از متفقین تقاضای بمباران خطوط ریلی به آشویتس را داشتند، "نجات‌طلب" بودند؟ آیا آرزواندیشی می‌کردند؟ یا واقع‌بینانه‌ترین ارزیابی ممکن را از وضعیت‌شان داشتند؟»\\
— مضمونی بازسازی‌شده از استدلال‌های منتقدان ضدمداخله‌گرایی مطلق
\end{tcolorbox}

این صدا را نمی‌توان با آمار «۸۴٪ مداخلات شکست خورده‌اند» ساکت کرد. زیرا برای کسی که در تاریکی مطلق زندگی می‌کند، حتی شمعی لرزان بهتر از هیچ است. پرسش واقعی این نیست که «آیا مداخله خطر دارد؟» — البته که دارد. پرسش واقعی این است: \textbf{«آیا هزینه‌ی مداخله، کمتر از هزینه‌ی عدم مداخله است؟»}

\begin{tcolorbox}[defbox, title={مفهوم کلیدی: نجات‌طلبی فداکارانه}]
\concept{نجات‌طلبی فداکارانه} (\lat{Sacrificial Liberation-Seeking}) رویکردی است که در آن یک ملت (یا بخشی از آن) آگاهانه و با چشم باز هزینه‌های مداخله‌ی خارجی — شامل خشونت، بی‌ثباتی، وابستگی موقت، و حتی تلفات نسلی — را می‌پذیرد، نه از سر آرزواندیشی یا خام‌اندیشی، بلکه بر اساس \concept{محاسبه‌ی عقلانی} مبنی بر اینکه:
\begin{enumerate}
\item وضع موجود (انسداد مطلق) بدتر از بدترین پیامد مداخله است
\item راه‌های داخلی تغییر واقعاً و نه فرضاً مسدود شده‌اند
\item هزینه‌ی پرداخت‌شده توسط نسل فعلی، سرمایه‌گذاری برای آزادی نسل‌های آینده است
\item جایگزین (عدم اقدام)، نه «صلح» بلکه \concept{خشونت ساختاری مداوم} است
\end{enumerate}
تفاوت بنیادین نجات‌طلبی فداکارانه با نجات‌طلبی ساده (بررسی‌شده در فصل‌های قبل): \textbf{آگاهی از هزینه‌ها + پذیرش آگاهانه + افق نسلی}.
\end{tcolorbox}


%% ================================================================
\section{پنج چارچوب پاسخ‌گویی: از کجا به مسئله نگاه کنیم؟}
\label{sec:ch15-frameworks}
%% ================================================================

\needspace{6\baselineskip}
پاسخ به پرسش «آیا نجات‌طلبی فداکارانه موجه است؟» عمیقاً به چارچوب تحلیلی‌ای بستگی دارد که از آن به مسئله نگاه می‌کنیم. پنج چارچوب اصلی وجود دارد:

\subsection{چارچوب اول: فایده‌گرایی نسلی}
\label{subsec:ch15-utilitarianism}

\concept{فایده‌گرایی نسلی} (\lat{Intergenerational Utilitarianism}) — مبتنی بر سنت فکری \thinker{جرمی بنتام} (\lat{Jeremy Bentham}) و \thinker{جان استوارت میل} (\lat{John Stuart Mill}) — اقدام درست را اقدامی می‌داند که بیشترین خیر (یا کمترین رنج) را برای بیشترین تعداد افراد، \textbf{در طول زمان}، تولید کند.\footnote{\lr{Parfit, Derek. \textit{Reasons and Persons}. Oxford: Clarendon Press, 1984, pp. 351–379.}}

\begin{tcolorbox}[wavebox, title={استدلال فایده‌گرایانه به نفع نجات‌طلبی فداکارانه}]
فرض کنید:
\begin{itemize}
\item یک دیکتاتوری مطلق سالانه ۱۰٬۰۰۰ نفر را می‌کشد و ده‌ها میلیون را در فقر و سرکوب نگه می‌دارد
\item هیچ راه داخلی‌ای برای تغییر وجود ندارد (سرکوب کامل)
\item مداخله‌ی خارجی ممکن است ۱۰۰٬۰۰۰ نفر تلفات فوری داشته باشد و ۱۰ سال بی‌ثباتی ایجاد کند
\item اما پس از آن، نظامی دموکراتیک (یا حداقل کمتر سرکوبگر) برقرار شود
\end{itemize}

\textbf{محاسبه:}
\begin{itemize}
\item هزینه‌ی تداوم وضع موجود (۵۰ سال): ۵۰۰٬۰۰۰ کشته + ده‌ها میلیون سرکوب‌شده
\item هزینه‌ی مداخله: ۱۰۰٬۰۰۰ تلفات + ۱۰ سال بی‌ثباتی
\item سود مداخله (اگر موفق شود): آزادی نسل‌های آینده
\end{itemize}

اگر \textbf{احتمال موفقیت} حداقل ۲۰–۳۰٪ باشد، فایده‌ی انتظاری (\lat{Expected Utility}) مداخله ممکن است از عدم مداخله بیشتر باشد.
\end{tcolorbox}

\thinker{درک پارفیت} (\lat{Derek Parfit}) در کتاب \textit{دلایل و اشخاص} استدلال کرده است که ما مسئولیت اخلاقی نسبت به نسل‌های آینده داریم — و این مسئولیت ممکن است ایجاب کند که نسل فعلی هزینه‌هایی را بپذیرد که به نفع آیندگان است.\footnote{\lr{Parfit, Derek. \textit{Reasons and Persons}, op. cit., pp. 351–379.}}

\subsection{چارچوب دوم: وظیفه‌گرایی کانتی}
\label{subsec:ch15-kantian}

\concept{وظیفه‌گرایی کانتی} (\lat{Kantian Deontology}) — مبتنی بر فلسفه‌ی \thinker{ایمانوئل کانت} (\lat{Immanuel Kant}) — بر حقوق بنیادین انسان تأکید دارد: هیچ‌کس نباید صرفاً ابزار اهداف دیگران باشد (\concept{امر مطلق} / \lat{Categorical Imperative}).

\begin{tcolorbox}[wavebox, title={استدلال کانتی: دو تفسیر متضاد}]
\textbf{تفسیر ضد مداخله (غالب در این کتاب):}
استفاده از ملت‌ها به‌عنوان «ابزار» اهداف ژئوپولیتیک مداخله‌گر، نقض امر مطلق کانتی است. مردم عراق «وسیله»ی اثبات قدرت آمریکا شدند — نه «غایت» فی‌نفسه.

\textbf{تفسیر حامی مداخله (بررسی‌شده در این فصل):}
تداوم دیکتاتوری نیز نقض امر مطلق است: دیکتاتور مردم خود را «ابزار» حفظ قدرتش می‌کند. اگر مداخله‌ی خارجی تنها راه بازگرداندن \concept{کرامت انسانی} (\lat{Human Dignity}) باشد، \textit{عدم} مداخله خود نقض وظیفه‌ی اخلاقی است.

\thinker{آلن بیوکنن} (\lat{Allen Buchanan})، فیلسوف سیاسی، استدلال کرده است که «حق حاکمیت» مشروط به رعایت حقوق اساسی شهروندان است — و دولتی که این حقوق را نقض کند، «حق» حاکمیت را از دست می‌دهد.\footnote{\lr{Buchanan, Allen. \textit{Justice, Legitimacy, and Self-Determination: Moral Foundations for International Law}. Oxford: Oxford University Press, 2004, pp. 233–270.}}
\end{tcolorbox}

\subsection{چارچوب سوم: عدالت بین‌نسلی}
\label{subsec:ch15-intergenerational}

\concept{عدالت بین‌نسلی} (\lat{Intergenerational Justice}) پرسش می‌کند: آیا نسل فعلی حق دارد — یا حتی \textit{وظیفه} دارد — برای آزادی نسل‌های آینده هزینه بپردازد؟

\thinker{جان رالز} (\lat{John Rawls}) در \textit{نظریه‌ی عدالت} مفهوم \concept{اصل پس‌انداز عادلانه} (\lat{Just Savings Principle}) را مطرح کرده: هر نسل وظیفه دارد شرایطی را فراهم کند که نسل‌های بعدی بتوانند از نهادهای عادلانه بهره‌مند شوند.\footnote{\lr{Rawls, John. \textit{A Theory of Justice}. Rev. ed., Cambridge, MA: Harvard University Press, 1999, pp. 251–258.}}

\begin{tcolorbox}[defbox, title={پرسش رالزی: پرده‌ی بی‌خبری بین‌نسلی}]
اگر شما پشت \concept{پرده‌ی بی‌خبری} (\lat{Veil of Ignorance}) رالزی بودید و نمی‌دانستید در کدام نسل به دنیا خواهید آمد — نسل فعلی (که هزینه‌ی مداخله را می‌پردازد) یا نسل آینده (که از آزادی بهره‌مند می‌شود) — چه تصمیمی می‌گرفتید؟

\textbf{استدلال حامیان مداخله:} شخص عقلانی پشت پرده‌ی بی‌خبری، مداخله‌ای را ترجیح می‌دهد که احتمال ۳۰٪ آزادی پایدار دارد، بر تداوم قطعی سرکوب مطلق — زیرا ریسک «زندگی در دوره‌ی بی‌ثباتی» کمتر از ریسک «زندگی در سرکوب ابدی» است.
\end{tcolorbox}

\subsection{چارچوب چهارم: واقع‌گرایی تراژیک}
\label{subsec:ch15-tragic-realism}

\concept{واقع‌گرایی تراژیک} (\lat{Tragic Realism}) — مبتنی بر آثار \thinker{رینهولد نیبور} (\lat{Reinhold Niebuhr}) و \thinker{مکس وبر} (\lat{Max Weber}) — می‌پذیرد که در سیاست، گاهی هیچ گزینه‌ی «خوبی» وجود ندارد — فقط گزینه‌های «بد» و «بدتر». وظیفه‌ی اخلاقی، انتخاب \concept{بد کمتر} (\lat{Lesser Evil}) است.\footnote{\lr{Niebuhr, Reinhold. \textit{Moral Man and Immoral Society: A Study in Ethics and Politics}. New York: Charles Scribner's Sons, 1932.}}

\begin{tcolorbox}[wavebox, title={واقع‌گرایی تراژیک و نجات‌طلبی فداکارانه}]
\thinker{مایکل والزر} در مقاله‌ی مهم «مشکل دست‌های آلوده» استدلال کرده است که رهبران سیاسی گاهی باید مرتکب «شر» شوند تا از «شر بزرگ‌تر» جلوگیری کنند — و این نه ضعف اخلاقی، بلکه مسئولیت اخلاقی است.\footnote{\lr{Walzer, Michael. "Political Action: The Problem of Dirty Hands," \textit{Philosophy \& Public Affairs}, Vol. 2, No. 2, 1973, pp. 160–180.}}

\thinker{مکس وبر} تمایز مشهوری میان \concept{اخلاق اعتقادی} (\lat{Ethic of Conviction}) و \concept{اخلاق مسئولیت} (\lat{Ethic of Responsibility}) قائل شد: اخلاق اعتقادی می‌گوید «مداخله‌ی خارجی بد است — نقطه»؛ اخلاق مسئولیت می‌پرسد «نتیجه‌ی عدم مداخله چیست؟ و آیا من مسئولیت آن نتیجه را می‌پذیرم؟»\footnote{\lr{Weber, Max. "Politics as a Vocation" (1919), in \textit{From Max Weber: Essays in Sociology}, trans. H.H. Gerth and C. Wright Mills. New York: Oxford University Press, 1946, pp. 77–128.}}
\end{tcolorbox}

\subsection{چارچوب پنجم: فلسفه‌ی مقاومت و رهایی}
\label{subsec:ch15-liberation-philosophy}

\concept{فلسفه‌ی رهایی} (\lat{Liberation Philosophy}) — مبتنی بر آثار \thinker{فرانتس فانون}، \thinker{پائولو فریره}، و \thinker{آنتونیو نگری} (\lat{Antonio Negri}) — تأکید می‌کند که ستمدیدگان نه‌فقط «حق» بلکه \concept{وظیفه} دارند علیه ستم بجنگند — و استفاده از هر ابزار ممکن (شامل کمک خارجی) مشروع است.

\thinker{فانون} در \textit{دوزخیان زمین} نوشت: «استعمارزده که تحت خشونت ساختاری زندگی می‌کند، حق دارد از خشونت برای رهایی استفاده کند — زیرا خشونت اول از سوی استعمارگر آمده.»\footnote{\lr{Fanon, Frantz. \textit{The Wretched of the Earth}, op. cit., pp. 1–62.}} اگرچه فانون بر عاملیت داخلی تأکید داشت، منطق استدلالش قابل تعمیم به کمک‌خواهی خارجی نیز هست: اگر خشونت ساختاری رژیم آنقدر شدید است که هیچ مقاومت داخلی‌ای ممکن نیست، آیا درخواست کمک خارجی نوعی «خود-دفاعی» نیست؟


%% ================================================================
\section{جدول تطبیقی پنج چارچوب}
\label{sec:ch15-frameworks-table}
%% ================================================================

\needspace{6\baselineskip}

\begin{table}[htbp]
\centering
\caption{پنج چارچوب پاسخ‌گویی به مسئله‌ی نجات‌طلبی فداکارانه}
\label{tab:ch15-frameworks}
\begin{adjustbox}{max width=\linewidth}
\begin{tabularx}{\linewidth}{
    >{\raggedleft\arraybackslash}p{2.5cm}
    >{\raggedleft\arraybackslash}p{2.5cm}
    >{\raggedleft\arraybackslash}p{3cm}
    >{\raggedleft\arraybackslash}X
    >{\centering\arraybackslash}p{1.5cm}
}
\toprule
\rowcolor{PrimaryDeep}
\textcolor{white}{\textbf{چارچوب}} &
\textcolor{white}{\textbf{متفکر اصلی}} &
\textcolor{white}{\textbf{معیار ارزیابی}} &
\textcolor{white}{\textbf{حکم درباره‌ی نجات‌طلبی فداکارانه}} &
\textcolor{white}{\textbf{موضع}} \\
\midrule
فایده‌گرایی نسلی & بنتام، میل، پارفیت & بیشترین خیر/کمترین رنج در طول زمان & مشروط: اگر فایده‌ی انتظاری مثبت باشد & مشروط \\
\rowcolor{NeutralLight}
وظیفه‌گرایی کانتی & کانت، بیوکنن & حقوق بنیادین و کرامت انسانی & دوگانه: هم له و هم علیه استدلال دارد & متناقض \\
عدالت بین‌نسلی & رالز، پارفیت & عدالت بین نسل‌ها & حمایت مشروط: وظیفه‌ی نسل فعلی نسبت به آیندگان & مشروط \\
\rowcolor{NeutralLight}
واقع‌گرایی تراژیک & نیبور، وبر، والزر & انتخاب بد کمتر & حمایت: وقتی همه‌ی گزینه‌ها بد هستند & حمایت \\
فلسفه‌ی رهایی & فانون، فریره & رهایی ستمدیدگان & حمایت مشروط: ابزار رهایی مشروع‌اند & حمایت \\
\bottomrule
\end{tabularx}
\end{adjustbox}
\end{table}


%% ================================================================
\section{شواهد تاریخی: مواردی که مداخله‌ی خارجی «نجات‌بخش» بود}
\label{sec:ch15-historical-evidence}
%% ================================================================

\needspace{6\baselineskip}
تحلیل منصفانه ایجاب می‌کند که به مواردی نگاه کنیم که مداخله‌ی خارجی — علی‌رغم هزینه‌های سنگین — \textit{واقعاً} به رهایی انجامید. فصل‌های قبلی عمدتاً بر شکست‌ها تمرکز داشتند؛ اکنون زمان بررسی موفقیت‌هاست — هرچند نسبی.

\subsection{سقوط نازیسم: مداخله‌ی پرهزینه‌ای که بدیلی نداشت}
\label{subsec:ch15-nazi}

\begin{tcolorbox}[casebox, title={مطالعه‌ی موردی: شکست آلمان نازی (۱۹۳۹–۱۹۴۵)}]
\textbf{واقعیت تلخ:} هیچ مقاومت داخلی‌ای در آلمان نازی نتوانست هیتلر را سرنگون کند. توطئه‌ی ۲۰ ژوئیه ۱۹۴۴ (عملیات والکوری / \lat{Operation Valkyrie}) — جدی‌ترین تلاش داخلی — شکست خورد.\footnote{\lr{Fest, Joachim C. \textit{Plotting Hitler's Death: The Story of the German Resistance}. New York: Metropolitan Books, 1996.}} مقاومت غیرنظامی (گروه رز سفید / \lat{White Rose}) الهام‌بخش اما بی‌اثر بود. دستگاه سرکوب نازی (گشتاپو، اس‌اس) آنقدر کارآمد بود که هرگونه مقاومت سازمان‌یافته را غیرممکن می‌ساخت.

\textbf{هزینه‌ی مداخله:} جنگ جهانی دوم حدود ۷۰–۸۵ میلیون کشته بر جای گذاشت — فجیع‌ترین فاجعه‌ی انسانی تاریخ.\footnote{\lr{Overy, Richard. "The Dictators: Hitler's Germany, Stalin's Russia," in \textit{The Oxford Companion to World War II}. Oxford: Oxford University Press, 2001.}}

\textbf{نتیجه:} آلمان غربی پس از ۱۹۴۵ به یکی از پایدارترین دموکراسی‌ها و قوی‌ترین اقتصادهای جهان تبدیل شد. ژاپن نیز همین مسیر را طی کرد.

\textbf{پرسش کلیدی:} آیا ۷۰ میلیون کشته «ارزشش» را داشت؟ این پرسش اخلاقاً غیرقابل‌پاسخ است. اما پرسش دیگری قابل پاسخ است: \textbf{آیا بدیلی وجود داشت؟} اگر متفقین مداخله نمی‌کردند، هولوکاست ادامه می‌یافت، ملت‌های اروپایی تحت اشغال می‌ماندند، و جهان بسیار تاریک‌تری شکل می‌گرفت.
\end{tcolorbox}

\thinker{تیموتی اسنایدر} (\lat{Timothy Snyder})، تاریخ‌دان ییل، در کتاب \textit{سرزمین‌های خونین} مستند کرده است که بین ۱۹۳۳ و ۱۹۴۵ حدود ۱۴ میلیون غیرنظامی — بیشترشان یهودی، لهستانی، اوکراینی، و بلاروسی — در «سرزمین‌های خونین» بین برلین و مسکو به قتل رسیدند.\footnote{\lr{Snyder, Timothy. \textit{Bloodlands: Europe Between Hitler and Stalin}. New York: Basic Books, 2010.}} عدم مداخله به معنای تداوم این کشتار بود.

\subsection{نجات اروپا از فاشیسم: طرح مارشال و دموکراسی‌سازی}
\label{subsec:ch15-marshall}

آنچه سقوط نازیسم را از سایر مداخلات (عراق، لیبی) متمایز می‌کند، فقط مداخله‌ی نظامی نبود — بلکه \concept{تعهد بلندمدت به بازسازی} بود:

\begin{tcolorbox}[successbox, title={عوامل موفقیت: بازسازی آلمان و ژاپن}]
\begin{enumerate}
\item \textbf{طرح مارشال:} ۱۳ میلیارد دلار (معادل ۱۵۰ میلیارد امروزی) برای بازسازی اروپا
\item \textbf{تعهد نسلی:} آمریکا هنوز هم — ۸۰ سال بعد — پایگاه نظامی در آلمان و ژاپن دارد
\item \textbf{بازسازی نهادی:} نه فقط بازسازی فیزیکی، بلکه بازسازی نهادهای دموکراتیک، قضایی، و آموزشی
\item \textbf{عدم انتقام‌جویی:} بر خلاف پیمان ورسای (۱۹۱۹) که آلمان را تحقیر کرد، رویکرد ۱۹۴۵ بازسازنده بود
\item \textbf{تهدید خارجی مشترک:} جنگ سرد اشغالگران و اشغال‌شدگان را متحد کرد
\item \textbf{سنت نهادی پیشین:} آلمان و ژاپن پیش از جنگ دولت‌های مدرن با بوروکراسی کارآمد بودند
\end{enumerate}
\end{tcolorbox}

\textbf{درس:} سقوط نازیسم نشان می‌دهد که مداخله‌ی خارجی \textit{می‌تواند} موفق باشد — اما تحت شرایط بسیار خاص و با تعهد نسلی. پرسش این است: آیا این شرایط در موارد دیگر هم وجود دارد؟

\subsection{شکست آپارتاید: ترکیب عاملیت داخلی و فشار خارجی}
\label{subsec:ch15-apartheid}

\begin{tcolorbox}[casebox, title={مطالعه‌ی موردی: پایان آپارتاید در آفریقای جنوبی (۱۹۹۴)}]
آفریقای جنوبی مثالی است از ترکیب موفق عاملیت داخلی و فشار خارجی:

\textbf{عامل داخلی:} \concept{کنگره‌ی ملی آفریقا} (\lat{ANC})، مقاومت مدنی گسترده، اعتصابات، و مبارزه‌ی مسلحانه‌ی محدود — \histref{نلسون ماندلا}{رهبر ANC} نماد مقاومت شد.\footnote{\lr{Mandela, Nelson. \textit{Long Walk to Freedom: The Autobiography of Nelson Mandela}. Boston: Little, Brown, 1994.}}

\textbf{فشار خارجی:}
\begin{itemize}
\item تحریم‌های بین‌المللی (ورزشی، اقتصادی، فرهنگی)
\item فشار افکار عمومی جهانی (جنبش ضد آپارتاید)
\item تحریم سرمایه‌گذاری (خروج شرکت‌های بزرگ)
\item فشار دیپلماتیک
\end{itemize}

\textbf{آنچه صورت نگرفت:} مداخله‌ی نظامی مستقیم.

\textbf{نتیجه:} انتقال مسالمت‌آمیز قدرت، دموکراسی، و \concept{کمیسیون حقیقت و آشتی} (\lat{Truth and Reconciliation Commission}).

\textbf{درس:} فشار خارجی مؤثرترین زمانی بود که \textit{ابزار} مبارزه‌ی داخلی شد — نه جایگزین آن. ماندلا از تحریم‌ها استقبال کرد اما هرگز «مداخله‌ی نظامی» نخواست.
\end{tcolorbox}

\subsection{رواندا ۱۹۹۴: قیمت عدم مداخله}
\label{subsec:ch15-rwanda}

همان‌گونه که در فصل~\ref{ch:contemporary} مفصلاً بحث شد، رواندا قوی‌ترین استدلال تاریخی به نفع مداخله است: ۸۰۰٬۰۰۰ نفر طی ۱۰۰ روز کشته شدند در حالی که جهان تماشا کرد.

ژنرال \histref{رومئو دالر}{فرمانده‌ی یونامیر} بعدها نوشت که با ۵٬۰۰۰ نیروی مجهز می‌توانست نسل‌کشی را متوقف کند.\footnote{\lr{Dallaire, Roméo. \textit{Shake Hands with the Devil}, op. cit., pp. 374–375.}} هزینه‌ی مداخله (احتمالاً چند صد کشته‌ی نظامی) در مقایسه با هزینه‌ی عدم مداخله (۸۰۰٬۰۰۰ کشته) ناچیز بود.

\begin{tcolorbox}[enemybox, title={رواندا: وقتی عدم مداخله جنایت است}]
رواندا نشان داد که \concept{عدم مداخله} نیز یک «تصمیم» است — تصمیمی با پیامدهای فاجعه‌بار. کسانی که می‌گویند «هرگز مداخله نکنید» باید پاسخگوی رواندا نیز باشند. \thinker{سامانتا پاور} به‌درستی اشاره کرده که «مشکلی از جهنم» هم عمل کردن و هم عمل نکردن است.\footnote{\lr{Power, Samantha. \textit{A Problem from Hell}, op. cit., p. 508.}}
\end{tcolorbox}

\subsection{کوزوو ۱۹۹۹: مداخله‌ی «غیرقانونی اما مشروع»}
\label{subsec:ch15-kosovo-success}

بمباران ناتو در کوزوو — بدون مجوز شورای امنیت — از نظر حقوقی مشکل‌ساز بود. اما از نظر نتایج:

\begin{itemize}
\item پاک‌سازی قومی متوقف شد
\item کوزوو در ۲۰۰۸ مستقل شد
\item اگرچه فقیر و وابسته، اما مردم کوزوو از سرکوب صربستان رها شدند
\item هیچ نسل‌کشی‌ای مانند سربرنیتسا (۱۹۹۵) تکرار نشد
\end{itemize}

آیا کوزووی‌ها باید منتظر «عاملیت جمعی» و «مقاومت مدنی غیرخشونت‌آمیز» می‌ماندند در حالی که میلوشویچ آن‌ها را قتل‌عام می‌کرد؟

\subsection{کامبوج ۱۹۷۹: مداخله‌ی ویتنام و پایان خمرهای سرخ}
\label{subsec:ch15-cambodia}

\begin{tcolorbox}[casebox, title={مطالعه‌ی موردی: کامبوج — مداخله‌ای که جان نجات داد}]
رژیم \histref{پل پوت}{رهبر خمرهای سرخ} بین ۱۹۷۵ و ۱۹۷۹ حدود ۱٫۵ تا ۲ میلیون نفر از ۸ میلیون جمعیت کامبوج را کشت. در ژانویه ۱۹۷۹، ویتنام به کامبوج حمله و رژیم خمرهای سرخ را سرنگون کرد.\footnote{\lr{Kiernan, Ben. \textit{The Pol Pot Regime: Race, Power, and Genocide in Cambodia under the Khmer Rouge, 1975–79}. 3rd ed., New Haven: Yale University Press, 2008.}}

\textbf{پارادوکس:}
\begin{itemize}
\item مداخله‌ی ویتنام احتمالاً صدها هزار جان نجات داد
\item اما جامعه‌ی بین‌المللی (آمریکا و چین) مداخله را محکوم کرد — زیرا ویتنام متحد شوروی بود
\item خمرهای سرخ تا سال‌ها بعد «نماینده‌ی قانونی» کامبوج در سازمان ملل باقی ماندند
\item غرب از نسل‌کشی پل پوت مطلع بود اما عمل نکرد — ویتنام عمل کرد
\end{itemize}

\textbf{درس:} مداخله‌ی «غیرمشروع» حقوقی (ویتنام بدون مجوز شورای امنیت حمله کرد) ممکن است اخلاقاً مشروع‌ترین اقدام ممکن باشد.
\end{tcolorbox}


%% ================================================================
\section{نظریه‌ها و استدلال‌های حمایت‌کننده}
\label{sec:ch15-theories}
%% ================================================================

\needspace{6\baselineskip}
\subsection{نظریه‌ی «جنگ عادلانه» و گسترش آن}
\label{subsec:ch15-just-war}

سنت \concept{جنگ عادلانه} (\lat{Just War Theory}) — از \thinker{آوگوستین قدیس} و \thinker{توماس آکویناس} تا \thinker{مایکل والزر} — شرایطی را تعریف می‌کند که تحت آن استفاده از زور اخلاقاً مجاز است. والزر در ویرایش جدید \textit{جنگ‌های عادلانه و ناعادلانه} استدلال کرده که مداخله‌ی انسان‌دوستانه تحت سه شرط مجاز است:\footnote{\lr{Walzer, Michael. \textit{Just and Unjust Wars}, op. cit., pp. 101–108.}}

\begin{enumerate}
\item نقض حقوق بشر آنقدر شدید باشد که «وجدان بشریت را تکان دهد» (\lat{shock the conscience of mankind})
\item مداخله‌ی نظامی تنها ابزار مؤثر برای توقف نقض باشد
\item نتیجه‌ی مداخله بهتر از عدم مداخله باشد (اصل \concept{تناسب} / \lat{proportionality})
\end{enumerate}

\subsection{نظریه‌ی «شکستن انسداد تاریخی»}
\label{subsec:ch15-historical-deadlock}

\thinker{بارینگتون مور} (\lat{Barrington Moore Jr.}) در کتاب کلاسیک \textit{ریشه‌های اجتماعی دیکتاتوری و دموکراسی} نشان داده است که گذار به دموکراسی تقریباً همیشه با خشونت همراه بوده — از انقلاب فرانسه و جنگ داخلی آمریکا تا جنگ جهانی دوم.\footnote{\lr{Moore, Barrington Jr. \textit{Social Origins of Dictatorship and Democracy: Lord and Peasant in the Making of the Modern World}. Boston: Beacon Press, 1966.}}

\begin{tcolorbox}[defbox, title={مفهوم: انسداد تاریخی}]
\concept{انسداد تاریخی} (\lat{Historical Deadlock}) وضعیتی است که در آن:
\begin{enumerate}
\item رژیم حاکم آنقدر قوی و سرکوبگر است که هیچ مقاومت داخلی‌ای امکان موفقیت ندارد
\item نخبگان رژیم هیچ انگیزه‌ای برای اصلاح ندارند (رانت کافی دارند)
\item جامعه‌ی مدنی مستقل نابود شده
\item فضای رسانه‌ای بسته است
\item ارتش وفادار به رژیم است
\item هیچ «شکافی» در ساختار قدرت وجود ندارد که از آن بتوان بهره برد
\end{enumerate}

در چنین وضعیتی، توصیه‌ی «مقاومت مدنی غیرخشونت‌آمیز» و «عاملیت جمعی» ممکن است \concept{آرزواندیشی معکوس} باشد: آرزواندیشی‌ای که به جای امید به بیرون، به «قدرت درون» امید می‌بندد — بدون توجه به واقعیت سرکوب مطلق.
\end{tcolorbox}

\subsection{استدلال «هزینه‌ی زمان»: وقتی انتظار خود خشونت است}
\label{subsec:ch15-cost-of-time}

\thinker{آمارتیا سن} (\lat{Amartya Sen})، اقتصاددان برنده‌ی نوبل، مفهوم \concept{خشونت ساختاری} (\lat{Structural Violence}) — مطرح‌شده توسط \thinker{یوهان گالتونگ} (\lat{Johan Galtung}) — را به حوزه‌ی اقتصاد توسعه وارد کرده است: فقر، سوءتغذیه، و بیماری‌های قابل‌پیشگیری خود اَشکال خشونت‌اند — حتی اگر کسی تفنگ شلیک نکند.\footnote{\lr{Sen, Amartya. \textit{Development as Freedom}. New York: Alfred A. Knopf, 1999.}}$^{,}$\footnote{\lr{Galtung, Johan. "Violence, Peace, and Peace Research," \textit{Journal of Peace Research}, Vol. 6, No. 3, 1969, pp. 167–191.}}

\begin{tcolorbox}[wavebox, title={استدلال هزینه‌ی زمان}]
هر سال که دیکتاتوری ادامه یابد:
\begin{itemize}
\item هزاران نفر کشته، زندانی، شکنجه، و تبعید می‌شوند
\item میلیون‌ها نفر در فقر و سرکوب زندگی می‌کنند
\item نخبگان و مغزها فرار می‌کنند → سرمایه‌ی انسانی نابود می‌شود
\item کودکان بدون آموزش کیفی بزرگ می‌شوند → نسل بعدی نیز تضعیف می‌شود
\item محیط زیست تخریب می‌شود (رژیم‌های دیکتاتوری معمولاً حساسیت زیست‌محیطی ندارند)
\end{itemize}

\textbf{«انتظار» خود نوعی تصمیم است — تصمیم به پذیرش خشونت ساختاری مداوم.}

اگر مداخله‌ی خارجی ۱۰ سال زودتر به دموکراسی برسد، ده سال خشونت ساختاری «پس‌انداز» شده — یعنی صدها هزار جان، میلیون‌ها سال آموزش، و تریلیون‌ها دلار سرمایه‌ی انسانی.
\end{tcolorbox}


%% ================================================================
\section{نمودار: مدل تصمیم‌گیری نجات‌طلبی فداکارانه}
\label{sec:ch15-model-tikz}
%% ================================================================

\needspace{14\baselineskip}

\begin{figure}[htbp]
\centering
\begin{tikzpicture}[
    >=Stealth,
    node distance=2cm,
    box/.style={rectangle, draw=PrimaryDeep, fill=PrimaryDeep!10, rounded corners=6pt,
        text width=3.5cm, minimum height=1.5cm, align=center, font=\small},
    cost/.style={rectangle, draw=red!70!black, fill=red!10, rounded corners=6pt,
        text width=3cm, minimum height=1cm, align=center, font=\footnotesize},
    benefit/.style={rectangle, draw=green!60!black, fill=green!10, rounded corners=6pt,
        text width=3cm, minimum height=1cm, align=center, font=\footnotesize},
    arrow/.style={->, thick, PrimaryDeep},
    scale=0.78, transform shape
]

% دو مسیر
\node[box] (start) at (0, 8) {\rl{انسداد سیاسی مطلق}\\[3pt]\rl{(«بالاتر از سیاهی...»)}};

% مسیر ۱: عدم اقدام
\node[box, draw=gray] (inaction) at (-5, 5) {\rl{مسیر ۱:}\\[3pt]\rl{عدم اقدام}\\[3pt]\rl{(«صبر»)}};
\node[cost] (cost1a) at (-7, 2) {\rl{خشونت ساختاری}\\[2pt]\rl{مداوم (سالانه)}};
\node[cost] (cost1b) at (-3, 2) {\rl{نابودی سرمایه‌ی}\\[2pt]\rl{انسانی نسل‌ها}};
\node[cost] (cost1c) at (-5, -0.5) {\rl{تداوم نامعلوم}\\[2pt]\rl{(۱۰ سال؟ ۵۰ سال؟ ابد؟)}};

% مسیر ۲: نجات‌طلبی فداکارانه
\node[box, draw=AccentGold] (action) at (5, 5) {\rl{مسیر ۲:}\\[3pt]\rl{نجات‌طلبی فداکارانه}\\[3pt]\rl{(مداخله + هزینه)}};
\node[cost] (cost2a) at (3, 2) {\rl{خشونت فوری:}\\[2pt]\rl{جنگ، بی‌ثباتی}};
\node[cost] (cost2b) at (7, 2) {\rl{وابستگی موقت}\\[2pt]\rl{به مداخله‌گر}};
\node[benefit] (ben2) at (5, -0.5) {\rl{احتمال (\%p):}\\[2pt]\rl{آزادی نسل‌های آینده}};

% فلش‌ها
\draw[arrow] (start) -- (inaction);
\draw[arrow] (start) -- (action);
\draw[arrow, red!70!black] (inaction) -- (cost1a);
\draw[arrow, red!70!black] (inaction) -- (cost1b);
\draw[arrow, red!70!black] (cost1a) -- (cost1c);
\draw[arrow, red!70!black] (cost1b) -- (cost1c);
\draw[arrow, red!70!black] (action) -- (cost2a);
\draw[arrow, red!70!black] (action) -- (cost2b);
\draw[arrow, green!60!black] (cost2a) -- (ben2);
\draw[arrow, green!60!black] (cost2b) -- (ben2);

% مقایسه
\node[rectangle, draw=PrimaryDeep, fill=yellow!15, rounded corners, font=\footnotesize,
    text width=8cm, align=center] (compare) at (0, -3)
    {\textbf{\rl{پرسش بنیادین:}}\\[4pt]
    \rl{آیا هزینه‌ی مسیر ۲ (جنگ + بی‌ثباتی + وابستگی) × احتمال موفقیت}\\[2pt]
    \rl{کمتر از هزینه‌ی مسیر ۱ (سرکوب ابدی + نابودی نسل‌ها) است؟}};

\draw[arrow, dashed, PrimaryDeep] (cost1c) -- (compare);
\draw[arrow, dashed, PrimaryDeep] (ben2) -- (compare);

% عنوان
\node[font=\normalsize\bfseries, PrimaryDeep] at (0, 10) {\rl{مدل تصمیم‌گیری: دو مسیر در انسداد مطلق}};

\end{tikzpicture}
\caption{مقایسه‌ی دو مسیر: عدم اقدام در برابر نجات‌طلبی فداکارانه}
\label{fig:ch15-decision-model}
\end{figure}


%% ================================================================
\section{بهترین استدلال‌ها به نفع نجات‌طلبی فداکارانه}
\label{sec:ch15-best-arguments}
%% ================================================================

\needspace{6\baselineskip}

\begin{tcolorbox}[successbox, title={ده استدلال قوی به نفع نجات‌طلبی فداکارانه}]

\textbf{استدلال ۱: «آلمان نازی»}
هیچ مقاومت داخلی‌ای نازیسم را شکست نداد. فقط مداخله‌ی خارجی توانست. اگر متفقین مداخله نمی‌کردند، هولوکاست ادامه می‌یافت. آیا ملت‌هایی که امروز در وضعیت مشابه‌اند باید منتظر «عاملیت جمعی» بمانند؟

\textbf{استدلال ۲: «هزینه‌ی فرصت نسلی»}
هر سال تأخیر در تغییر = یک نسل کودک محروم از آموزش کیفی، یک نسل جوان محروم از فرصت، و یک نسل نخبه‌ی فراری. هزینه‌ی «صبر» صفر نیست — بسیار بالاست.

\textbf{استدلال ۳: «رواندا معکوس»}
اگر مداخله نشود و فاجعه رخ دهد، چه کسی مسئول است؟ «عدم مداخله» نیز یک تصمیم اخلاقی است — و گاهی تصمیمی جنایتکارانه.

\textbf{استدلال ۴: «مقاومت مدنی همیشه ممکن نیست»}
پژوهش چنووت و استفان (که این کتاب بارها به آن استناد کرد) نشان می‌دهد مقاومت مدنی «بیشتر» موفق است — نه «همیشه». در رژیم‌های توتالیتر با سرکوب کامل (کره‌ی شمالی، ترکمنستان، اریتره)، مقاومت مدنی عملاً غیرممکن است.

\textbf{استدلال ۵: «مسئولیت نسبت به آیندگان»}
\thinker{ادموند برک} (\lat{Edmund Burke}) جامعه را «قراردادی میان مردگان، زندگان، و آنان که هنوز زاده نشده‌اند» نامید.\footnote{\lr{Burke, Edmund. \textit{Reflections on the Revolution in France}. London: J. Dodsley, 1790.}} اگر نسل فعلی وظیفه‌ای نسبت به آیندگان دارد، آیا این وظیفه شامل «شکستن زنجیر» نیز نمی‌شود — حتی به هزینه‌ی خود؟

\textbf{استدلال ۶: «سوگیری وضع موجود»}
\concept{سوگیری وضع موجود} (\lat{Status Quo Bias}) — بررسی‌شده در فصل~\ref{ch:social-psychology} — نشان می‌دهد که انسان‌ها ریسک تغییر را بیش از ریسک تداوم وضع موجود ارزیابی می‌کنند. شاید مخالفت با نجات‌طلبی نیز تا حدی ناشی از این سوگیری شناختی باشد — ترس از ناشناخته‌ی تغییر، بیش از ترس از شناخته‌شده‌ی سرکوب.

\textbf{استدلال ۷: «اصل بردگی»}
هیچ‌کس امروز استدلال نمی‌کند که بردگان آمریکایی باید «صبر می‌کردند تا جامعه‌ی سفیدپوست خودش به اشتباهش پی ببرد.» جنگ داخلی آمریکا (۱۸۶۱–۱۸۶۵) با ۶۲۰٬۰۰۰ کشته پایان یافت — اما بردگی را لغو کرد. آیا کسی می‌تواند بگوید «ارزشش را نداشت»؟ بردگان فراری که به ارتش اتحادیه پیوستند، «نجات‌طلب» نبودند — رهایی‌جو بودند.\footnote{\lr{McPherson, James M. \textit{Battle Cry of Freedom: The Civil War Era}. New York: Oxford University Press, 1988, pp. 355–358.}}

\textbf{استدلال ۸: «مغالطه‌ی صلح منفی»}
\thinker{مارتین لوتر کینگ} در \textit{نامه از زندان بیرمنگام} تمایز مهمی میان \concept{صلح منفی} (غیاب تنش) و \concept{صلح مثبت} (حضور عدالت) قائل شد.\footnote{\lr{King, Martin Luther Jr. "Letter from Birmingham Jail," 1963, in \textit{Why We Can't Wait}. New York: Harper \& Row, 1964.}} دیکتاتوری ممکن است «صلح» (= غیاب جنگ آشکار) داشته باشد — اما این «صلح» استوار بر سرکوب، شکنجه، و خشونت ساختاری است. مخالفت با مداخله به‌نام «حفظ صلح» در واقع حفظ خشونت پنهان است.

\textbf{استدلال ۹: «اصل نجات غریق»}
\thinker{پیتر سینگر} (\lat{Peter Singer}) استدلال مشهوری دارد: اگر از کنار استخری رد شوید و کودکی را در حال غرق‌شدن ببینید، وظیفه‌ی اخلاقی شما نجات اوست — حتی اگر لباس‌تان خیس شود.\footnote{\lr{Singer, Peter. "Famine, Affluence, and Morality," \textit{Philosophy \& Public Affairs}, Vol. 1, No. 3, 1972, pp. 229–243.}} آیا این وظیفه در مقیاس ملت‌ها نیز صادق نیست؟ اگر ملتی در حال «غرق‌شدن» در نسل‌کشی یا سرکوب مطلق است، آیا جامعه‌ی بین‌المللی وظیفه‌ی نجات ندارد — و آیا خود آن ملت حق ندارد فریاد کمک بزند؟

\textbf{استدلال ۱۰: «واقع‌گرایی محض»}
فارغ از تمام ملاحظات اخلاقی، یک واقعیت ساده وجود دارد: بعضی رژیم‌ها از درون تغییرناپذیرند. کره‌ی شمالی ۷۵ سال است که دیکتاتوری مطلق دارد. اریتره ۳۰ سال. ترکمنستان ۳۳ سال. هیچ نشانه‌ای از تغییر داخلی در هیچ‌یک از این‌ها دیده نمی‌شود. آیا توصیه‌ی «صبر کنید، عاملیت جمعی خودش کار می‌کند» به مردم کره‌ی شمالی، واقع‌بینانه است — یا مصداق بارز همان «آرزواندیشی» است که این کتاب از آن انتقاد می‌کند؟
\end{tcolorbox}


%% ================================================================
\section{پاسخ مدافع: چرا این استدلال‌ها، با همه‌ی قوتشان، کافی نیستند}
\label{sec:ch15-counterarguments}
%% ================================================================

\needspace{6\baselineskip}
اکنون نوبت آن است که نویسنده از نقش «مدافع شیطان» خارج شود و به ارزیابی انتقادی این استدلال‌ها بپردازد. هدف نه رد کامل آن‌ها — که صادقانه بگوییم بعضی‌شان بسیار قوی‌اند — بلکه نشان‌دادن \concept{نقاط کور} و \concept{مفروضات پنهان} هر استدلال است.

\subsection{نقد اول: مغالطه‌ی تعمیم از استثنا}
\label{subsec:ch15-critique-generalization}

\begin{tcolorbox}[enemybox, title={مشکل بنیادین: آلمان نازی استثناست، نه قاعده}]
قوی‌ترین استدلال حامیان نجات‌طلبی فداکارانه — مورد آلمان نازی — دقیقاً به همین دلیل قوی است که \textbf{استثنایی} است. شرایط آلمان ۱۹۳۹–۱۹۴۵ تکرارناپذیر بود:

\begin{itemize}
\item آلمان نازی نه فقط مردم خود، بلکه \textit{کل اروپا} را تهدید می‌کرد — مداخله‌ی متفقین اساساً \concept{دفاع از خود} بود، نه «مداخله‌ی انسان‌دوستانه»
\item هزینه‌ی عدم مداخله برای \textit{خود متفقین} نجومی بود (اشغال بریتانیا، سلطه‌ی نازی بر اروپا)
\item پس از جنگ، منافع ژئوپولیتیک متفقین (مهار شوروی) با بازسازی آلمان و ژاپن \textit{هم‌راستا} بود
\item آلمان و ژاپن پیش از جنگ جوامع صنعتی مدرن با بوروکراسی کارآمد بودند
\end{itemize}

\textbf{هیچ‌یک از این شرایط} در مورد عراق، لیبی، افغانستان، سوریه، یا اکثر کشورهای امروزی صادق نیست. تعمیم از آلمان نازی به وضعیت‌های معاصر \concept{مغالطه‌ی ترکیب} (\lat{Fallacy of Composition}) است.\footnote{\lr{Pei, Minxin and Sara Kasper. "Lessons from the Past: The American Record on Nation Building," \textit{Carnegie Endowment Policy Brief}, No. 24, 2003.}}
\end{tcolorbox}

\thinker{مینکسین پِی} (\lat{Minxin Pei}) و \thinker{سارا کاسپر} در پژوهش جامعی برای بنیاد کارنگی، ۱۶ مورد «ملت‌سازی» آمریکایی از ۱۹۰۰ تا ۲۰۰۳ را بررسی کردند و نتیجه گرفتند که فقط در ۴ مورد (آلمان، ژاپن، گرنادا، پاناما) دموکراسی پایدار شکل گرفت — و هر چهار مورد شرایط بسیار خاصی داشتند.\footnote{\lr{Pei and Kasper, op. cit., pp. 3–8.}} نرخ موفقیت: \textbf{۲۵٪}. و حتی این ۲۵٪ هم با شرایطی بسیار خاص.

\subsection{نقد دوم: مسئله‌ی «چه کسی تصمیم می‌گیرد؟»}
\label{subsec:ch15-critique-who-decides}

\begin{tcolorbox}[wavebox, title={پارادوکس نمایندگی}]
نجات‌طلبی فداکارانه فرض می‌کند که «ملت» آگاهانه هزینه‌ی مداخله را می‌پذیرد. اما:

\begin{enumerate}
\item \textbf{چه کسی به نمایندگی از «ملت» تصمیم می‌گیرد؟} در دیکتاتوری، انتخابات آزاد وجود ندارد. رفراندوم ممکن نیست. چگونه می‌توان مطمئن شد که «مردم» مداخله می‌خواهند — و نه فقط بخشی از اپوزیسیون تبعیدی؟
\item \textbf{آیا نسل فعلی حق دارد به جای نسل آینده تصمیم بگیرد؟} نجات‌طلبی فداکارانه می‌گوید «ما هزینه می‌دهیم تا آیندگان آزاد شوند.» اما نسل‌های آینده ممکن بود مسیر دیگری ترجیح دهند. این تصمیم‌گیری به جای دیگران، خود نوعی \concept{پدرسالاری} (\lat{Paternalism}) است.
\item \textbf{مسئله‌ی اطلاعات:} چگونه می‌توان «آگاهانه» تصمیم گرفت وقتی اطلاعات ناقص است؟ مردم عراق نمی‌دانستند مداخله‌ی آمریکا چه عواقبی خواهد داشت. مردم لیبی نمی‌دانستند پس از قذافی چه خواهد شد. «آگاهی از هزینه» در تئوری ساده است — در عمل، غیرممکن.
\end{enumerate}
\end{tcolorbox}

\thinker{هانا آرنت} (\lat{Hannah Arendt}) هشدار داده بود که خشونت ذاتاً \concept{غیرقابل‌پیش‌بینی} است: «هیچ‌کس نمی‌تواند نتایج خشونت را پیش‌بینی کند — خشونت منطق خودش را دارد و از کنترل آغازگرانش خارج می‌شود.»\footnote{\lr{Arendt, Hannah. \textit{On Violence}. New York: Harcourt, 1970, pp. 4–8.}} اگر خشونت غیرقابل‌پیش‌بینی است، «محاسبه‌ی عقلانی هزینه‌ها» خود نوعی توهم است.

\subsection{نقد سوم: مسئله‌ی انگیزه‌ی مداخله‌گر}
\label{subsec:ch15-critique-motives}

\begin{tcolorbox}[defbox, title={واقعیت تلخ: هیچ مداخله‌ی «خیرخواهانه‌ی خالص» وجود ندارد}]
تمام مداخلات تاریخی مبتنی بر \concept{منافع استراتژیک} مداخله‌گر بوده‌اند — حتی «موفق»ترین آن‌ها:

\begin{itemize}
\item \textbf{جنگ جهانی دوم:} آمریکا تا حمله‌ی ژاپن به پرل هاربر (دسامبر ۱۹۴۱) وارد جنگ نشد — نه به‌خاطر هولوکاست، بلکه به‌خاطر تهدید مستقیم
\item \textbf{کوزوو ۱۹۹۹:} ناتو عمل کرد — اما در رواندا ۱۹۹۴ عمل نکرد. تفاوت: رواندا در «حیاط خلوت» اروپا نبود
\item \textbf{کامبوج ۱۹۷۹:} ویتنام به‌خاطر تنش‌های مرزی و رقابت منطقه‌ای حمله کرد — نه صرفاً برای نجات کامبوجی‌ها
\item \textbf{لیبی ۲۰۱۱:} فرانسه و بریتانیا برای تغییر رژیم و منافع نفتی مداخله کردند — نه «مسئولیت حمایت»
\end{itemize}

\textbf{نتیجه:} نجات‌طلب فداکارانه ممکن است «آگاهانه» هزینه بپذیرد — اما مداخله‌گر خارجی منافع خودش را دنبال می‌کند. «ابزار» رهایی (مداخله‌گر) اهداف متفاوتی از «غایت» رهایی (ملت ستمدیده) دارد. این \concept{عدم تقارن انگیزه‌ها} (\lat{Incentive Asymmetry}) ریشه‌ی اصلی شکست اکثر مداخلات است.\footnote{\lr{Mearsheimer, John J. \textit{The Great Delusion: Liberal Dreams and International Realities}. New Haven: Yale University Press, 2018, pp. 150–189.}}
\end{tcolorbox}

\subsection{نقد چهارم: مسئله‌ی «بالاتر از سیاهی» واقعاً هست}
\label{subsec:ch15-critique-worse-than-black}

عنوان این فصل — «بالاتر از سیاهی رنگی نیست» — فرض می‌کند که وضعیت نمی‌تواند بدتر شود. اما تاریخ نشان می‌دهد که \textbf{بالاتر از سیاهی، سیاهی عمیق‌تری وجود دارد}:

\begin{tcolorbox}[enemybox, title={وقتی مداخله وضع را بدتر کرد: سه نمونه}]
\textbf{عراق ۲۰۰۳:}
\begin{itemize}
\item قبل از مداخله: دیکتاتوری صدام — سرکوبگر اما «باثبات»
\item پس از مداخله: ۲۰۰٬۰۰۰–۶۰۰٬۰۰۰ کشته، جنگ فرقه‌ای، ظهور داعش، میلیون‌ها آواره
\item ۲۰ سال بعد: عراق هنوز بی‌ثبات، فاسد، و تحت نفوذ خارجی
\end{itemize}

\textbf{لیبی ۲۰۱۱:}
\begin{itemize}
\item قبل از مداخله: دیکتاتوری قذافی — عجیب اما نسبتاً باثبات
\item پس از مداخله: دو دولت رقیب، جنگ داخلی مداوم، بازار برده‌فروشی، پایگاه ترور
\item ۱۳ سال بعد: لیبی هنوز «دولت شکست‌خورده» است
\end{itemize}

\textbf{افغانستان ۲۰۰۱–۲۰۲۱:}
\begin{itemize}
\item ۲۰ سال اشغال، ۲ تریلیون دلار هزینه، هزاران کشته
\item نتیجه: بازگشت طالبان به قدرت در ۲۰۲۱ — دقیقاً وضعیت ۲۰۰۱
\item زنان افغانستان امروز وضعیت بدتری نسبت به ۲۰۰۱ دارند — زیرا طالبان انتقام‌جوتر بازگشته
\end{itemize}
\end{tcolorbox}

\textbf{درس:} فرض «بالاتر از سیاهی رنگی نیست» اغلب نادرست است. مداخله‌ی ناموفق می‌تواند وضعیت را از «سیاه» به «سیاه‌تر» ببرد — جنگ داخلی، تجزیه، ظهور گروه‌های تروریستی، و فروپاشی کامل نظم اجتماعی.


\subsection{نقد پنجم: مسئله‌ی فایده‌گرایی و ارقام}
\label{subsec:ch15-critique-utilitarian-numbers}

محاسبه‌ی فایده‌گرایانه‌ای که در بخش~\ref{subsec:ch15-utilitarianism} ارائه شد، مشکلات جدی دارد:

\begin{tcolorbox}[wavebox, title={پنج مشکل با محاسبه‌ی فایده‌گرایانه}]
\begin{enumerate}
\item \textbf{مشکل احتمال:} «احتمال موفقیت ۲۰–۳۰٪» از کجا می‌آید؟ پیش از مداخله، هیچ‌کس نمی‌داند احتمال واقعی چقدر است. پژوهش‌ها نشان می‌دهند انسان‌ها به‌طور سیستماتیک احتمال موفقیت اقدامات پرریسک را \concept{بیش‌برآورد} می‌کنند (\concept{خوش‌بینی غیرواقعی} / \lat{Optimism Bias}).\footnote{\lr{Kahneman, Daniel. \textit{Thinking, Fast and Slow}. New York: Farrar, Straus and Giroux, 2011, pp. 249–254.}}

\item \textbf{مشکل شمارش:} آیا «تعداد کشته‌ها» تنها معیار است؟ آیا تروما، آوارگی، تخریب فرهنگی، از بین رفتن اعتماد اجتماعی، و فروپاشی نهادها «قابل شمارش» هستند؟

\item \textbf{مشکل مقایسه‌ناپذیری:} آیا مرگ ۱۰۰٬۰۰۰ نفر در جنگ «معادل» مرگ ۱۰۰٬۰۰۰ نفر از فقر و بیماری در طول ۱۰ سال است؟ \thinker{برنارد ویلیامز} (\lat{Bernard Williams}) استدلال کرده که فایده‌گرایی با «یکسان‌سازی» انواع متفاوت رنج، عمق اخلاقی مسئله را از دست می‌دهد.\footnote{\lr{Williams, Bernard. "A Critique of Utilitarianism," in J.J.C. Smart and B. Williams, \textit{Utilitarianism: For and Against}. Cambridge: Cambridge University Press, 1973, pp. 77–150.}}

\item \textbf{مشکل نرخ تنزیل:} فایده‌گرایی نسلی فرض می‌کند رنج نسل آینده «ارزش» مساوی با رنج نسل فعلی دارد. اما آیا واقعاً می‌توان رنج انسان‌هایی که هنوز متولد نشده‌اند را با رنج انسان‌های زنده مقایسه کرد؟

\item \textbf{مشکل دامنه:} محاسبه فرض می‌کند «پس از مداخله، دموکراسی برقرار می‌شود.» اما اگر نشود چه؟ اگر مداخله به جنگ داخلی ۲۰ ساله منجر شود (مانند سوریه)؟ در آن صورت، تمام محاسبه ویران می‌شود.
\end{enumerate}
\end{tcolorbox}


\subsection{نقد ششم: خطر ابزاری‌شدن نجات‌طلبی فداکارانه}
\label{subsec:ch15-critique-instrumentalization}

شاید خطرناک‌ترین جنبه‌ی نجات‌طلبی فداکارانه، قابلیت \concept{سوءاستفاده} از آن است:

\begin{tcolorbox}[enemybox, title={چگونه قدرت‌ها از «نجات‌طلبی فداکارانه» سوءاستفاده می‌کنند}]
\begin{enumerate}
\item \textbf{تولید «تقاضا»ی مداخله:} قدرت‌های خارجی می‌توانند از طریق تأمین مالی گروه‌های اپوزیسیون، رسانه‌ها، و سازمان‌های غیردولتی، «تقاضا»ی مداخله را \textit{تولید} کنند — حتی اگر اکثریت مردم آن را نخواهند. تفکیک «تقاضای واقعی مردم» از «تقاضای مهندسی‌شده» بسیار دشوار است.

\item \textbf{مشروعیت‌بخشی به تجاوز:} هر قدرتی می‌تواند مداخله‌ی خود را «فداکارانه» و «برای نجات مردم» معرفی کند. روسیه حمله به اوکراین را «آزادسازی» نامید. آمریکا حمله به عراق را «آزادی عراق» (\lat{Operation Iraqi Freedom}) نامید. اگر «نجات‌طلبی فداکارانه» به‌عنوان اصل اخلاقی پذیرفته شود، هر مداخله‌ای می‌تواند زیر این چتر مشروعیت یابد.

\item \textbf{تضعیف حق حاکمیت:} پذیرش اصل «حق مداخله در صورت انسداد مطلق» عملاً حق حاکمیت ملی را زیر سؤال می‌برد — و این «حق» معمولاً فقط علیه کشورهای ضعیف اعمال می‌شود. هیچ‌کس پیشنهاد «مداخله‌ی انسان‌دوستانه» در چین (برای اویغورها) یا روسیه (برای چچن‌ها) نمی‌دهد.
\end{enumerate}
\end{tcolorbox}


%% ================================================================
\section{جدول تطبیقی: موفقیت‌ها و شکست‌ها — چه چیزی تفاوت ایجاد می‌کند؟}
\label{sec:ch15-comparison-table}
%% ================================================================

\needspace{6\baselineskip}

\begin{table}[htbp]
\centering
\caption{مقایسه‌ی شرایط مداخلات «موفق» و «شکست‌خورده»}
\label{tab:ch15-success-failure}
\begin{adjustbox}{max width=\linewidth}
\begin{tabularx}{\linewidth}{
    >{\raggedleft\arraybackslash}p{3cm}
    >{\raggedleft\arraybackslash}X
    >{\raggedleft\arraybackslash}X
}
\toprule
\rowcolor{PrimaryDeep}
\textcolor{white}{\textbf{عامل}} &
\textcolor{white}{\textbf{مداخلات «موفق»}} &
\textcolor{white}{\textbf{مداخلات «شکست‌خورده»}} \\
\midrule
سنت دولت‌سازی & دولت مدرن پیشین (آلمان، ژاپن) & فقدان دولت مرکزی (افغانستان، لیبی، سومالی) \\
\rowcolor{NeutralLight}
تعهد بازسازی & بلندمدت و نسلی (طرح مارشال) & کوتاه‌مدت و نتیجه‌گرا («مأموریت انجام شد») \\
همگنی جامعه & نسبتاً همگن (آلمان، ژاپن) & شکاف‌های عمیق قومی/مذهبی (عراق، لیبی، سوریه) \\
\rowcolor{NeutralLight}
منافع مداخله‌گر & هم‌راستا با بازسازی (جنگ سرد) & متعارض با بازسازی (نفت، ژئوپولیتیک) \\
رویکرد پس از جنگ & بازسازنده، غیرانتقام‌جویانه & انحلال نهادها (بعث‌زدایی در عراق) \\
\rowcolor{NeutralLight}
حمایت بین‌المللی & گسترده (سازمان ملل، متفقین) & محدود و اختلافی (تقسیم شورای امنیت) \\
سرمایه‌گذاری اقتصادی & عظیم (طرح مارشال) & ناکافی یا نامتوازن \\
\rowcolor{NeutralLight}
مدت‌زمان تعهد & دهه‌ها (آلمان: ۸۰+ سال) & چند سال (لیبی: خروج فوری پس از سقوط قذافی) \\
\bottomrule
\end{tabularx}
\end{adjustbox}
\end{table}

این جدول نکته‌ای حیاتی را آشکار می‌کند: \textbf{تفاوت میان موفقیت و شکست مداخله، در خود «مداخله» نیست — بلکه در شرایط پیش از مداخله و تعهد پس از مداخله است.} مداخله بدون زمینه‌ی مناسب و بدون تعهد بلندمدت، تقریباً محکوم به شکست است.


%% ================================================================
\section{ترکیب: هفت شرط لازم برای نجات‌طلبی فداکارانه‌ی موجه}
\label{sec:ch15-conditions}
%% ================================================================

\needspace{6\baselineskip}
بر اساس تحلیل‌های این فصل — هم استدلال‌های مؤید و هم انتقادات — می‌توان هفت شرط لازم (نه کافی) برای نجات‌طلبی فداکارانه‌ی «موجه» استخراج کرد:

\begin{tcolorbox}[defbox, title={هفت شرط لازم برای نجات‌طلبی فداکارانه‌ی موجه}]
\begin{enumerate}
\item \textbf{انسداد مطلق اثبات‌شده:} باید نشان داده شود — نه فقط ادعا — که تمام راه‌های داخلی تغییر (مقاومت مدنی، اصلاحات تدریجی، فشار نخبگان، تغییر درون‌رژیمی) واقعاً و نه فرضاً آزموده شده و شکست خورده‌اند. \concept{بار اثبات} (\lat{Burden of Proof}) بر عهده‌ی مدعیان نجات‌طلبی است.

\item \textbf{تقاضای مردمی واقعی (نه مهندسی‌شده):} باید شواهد قوی وجود داشته باشد که \textit{اکثریت} مردم (نه فقط گروهی از تبعیدیان یا اپوزیسیون تأمین‌مالی‌شده) خواهان مداخله هستند — و از هزینه‌های آن آگاه‌اند.

\item \textbf{مداخله‌گر قابل‌اعتماد:} مداخله‌گر باید ثابت کرده باشد که توانایی و \textit{اراده‌ی} تعهد بلندمدت را دارد. مداخله‌ی «بزن‌و‌در‌رو» (\lat{Hit-and-Run Intervention}) نه فقط بی‌فایده، بلکه مخرب‌تر از عدم مداخله است.

\item \textbf{نقشه‌ی پس از سقوط:} باید برنامه‌ی مشخص، واقع‌بینانه، و تأمین‌مالی‌شده‌ای برای دوره‌ی پس از سقوط رژیم وجود داشته باشد — شامل بازسازی نهادی، عدالت انتقالی، و توسعه‌ی اقتصادی.

\item \textbf{حمایت بین‌المللی گسترده:} مداخله‌ی یک‌جانبه (مانند عراق ۲۰۰۳) بسیار بیشتر از مداخله‌ی چندجانبه (مانند کره‌ی جنوبی ۱۹۵۰) مستعد شکست است.

\item \textbf{تناسب:} هزینه‌ی مورد انتظار مداخله باید — بر اساس بهترین تخمین‌های موجود — کمتر از هزینه‌ی مورد انتظار عدم مداخله باشد. این تخمین باید \concept{محافظه‌کارانه} (\lat{Conservative}) باشد — یعنی بدترین سناریوی مداخله با بهترین سناریوی عدم مداخله مقایسه شود (اصل \concept{حداکثر بدبینی} / \lat{Maximin}).

\item \textbf{سازوکار خروج و انتقال:} از همان ابتدا باید مشخص باشد مداخله‌گر چه زمانی و چگونه خارج می‌شود و قدرت را به نهادهای بومی منتقل می‌کند. مداخله‌ی بدون استراتژی خروج = اشغال.
\end{enumerate}
\end{tcolorbox}

\begin{tcolorbox}[wavebox, title={آزمون عملی: کدام مداخلات تاریخی هفت شرط را داشتند؟}]
\begin{itemize}
\item \textbf{آلمان/ژاپن ۱۹۴۵:} تقریباً تمام هفت شرط برآورده بود → \textit{موفقیت}
\item \textbf{کره‌ی جنوبی ۱۹۵۰:} اکثر شرایط (به‌ویژه حمایت بین‌المللی و تعهد بلندمدت) → \textit{موفقیت نسبی}
\item \textbf{کوزوو ۱۹۹۹:} بیشتر شرایط (به‌جز حمایت کامل بین‌المللی) → \textit{موفقیت نسبی}
\item \textbf{عراق ۲۰۰۳:} حداکثر ۲ شرط از ۷ → \textit{شکست فاجعه‌بار}
\item \textbf{لیبی ۲۰۱۱:} حداکثر ۱ شرط از ۷ (فقط تناسب اولیه) → \textit{شکست}
\item \textbf{افغانستان ۲۰۰۱:} ۳ شرط در ابتدا، اما تعهد بلندمدت و نقشه‌ی بازسازی ناکافی → \textit{شکست}
\item \textbf{رواندا ۱۹۹۴ (مداخله‌ی فرضی):} اگر صورت می‌گرفت، اکثر شرایط فراهم بود → \textit{احتمالاً موفق}
\end{itemize}

\textbf{الگو:} هرچه تعداد شرایط برآورده‌شده بیشتر، احتمال موفقیت بیشتر. اما حتی برآوردن تمام شرایط نیز \textit{تضمین} موفقیت نمی‌کند — فقط احتمالش را افزایش می‌دهد.
\end{tcolorbox}


%% ================================================================
\section{نمودار: طیف مداخله — از فشار نرم تا مداخله‌ی نظامی}
\label{sec:ch15-intervention-spectrum}
%% ================================================================

\needspace{12\baselineskip}
نجات‌طلبی فداکارانه لزوماً به معنای مداخله‌ی نظامی نیست. طیف وسیعی از اشکال فشار خارجی وجود دارد:

\begin{figure}[htbp]
\centering
\begin{tikzpicture}[
    >=Stealth,
    scale=0.82, transform shape,
    level/.style={rectangle, draw=PrimaryDeep, rounded corners=5pt,
        minimum width=3cm, minimum height=1.2cm, align=center, font=\small},
]

% محور
\draw[->, ultra thick, PrimaryDeep!60] (-7, 0) -- (7, 0);
\node[font=\footnotesize, below] at (-6.5, -0.3) {\rl{هزینه‌ی کم / ریسک کم}};
\node[font=\footnotesize, below] at (6.5, -0.3) {\rl{هزینه‌ی بالا / ریسک بالا}};

% سطوح
\node[level, fill=green!15] (l1) at (-5.5, 1.5) {\rl{فشار دیپلماتیک}\\[2pt]\rl{و بیانیه‌ها}};
\node[level, fill=green!10] (l2) at (-2.5, 1.5) {\rl{تحریم‌های}\\[2pt]\rl{هدفمند}};
\node[level, fill=yellow!15] (l3) at (0.5, 1.5) {\rl{حمایت از}\\[2pt]\rl{جامعه‌ی مدنی}};
\node[level, fill=orange!15] (l4) at (3.5, 1.5) {\rl{منطقه‌ی}\\[2pt]\rl{پرواز ممنوع}};
\node[level, fill=red!15] (l5) at (6, 1.5) {\rl{مداخله‌ی}\\[2pt]\rl{نظامی مستقیم}};

% نتایج
\node[font=\scriptsize, text width=2.5cm, align=center] at (-5.5, 3.2)
    {\rl{مثال: آپارتاید}\\[2pt]\rl{(مؤثر با تداوم)}};
\node[font=\scriptsize, text width=2.5cm, align=center] at (-2.5, 3.2)
    {\rl{مثال: ایران}\\[2pt]\rl{(نتایج مبهم)}};
\node[font=\scriptsize, text width=2.5cm, align=center] at (0.5, 3.2)
    {\rl{مثال: اروپای شرقی}\\[2pt]\rl{(مؤثر در ترکیب)}};
\node[font=\scriptsize, text width=2.5cm, align=center] at (3.5, 3.2)
    {\rl{مثال: لیبی ۲۰۱۱}\\[2pt]\rl{(لغزنده‌ی شیب)}};
\node[font=\scriptsize, text width=2.5cm, align=center] at (6, 3.2)
    {\rl{مثال: عراق ۲۰۰۳}\\[2pt]\rl{(ریسک فاجعه)}};

% فلش‌ها
\foreach \x in {l1, l2, l3, l4, l5}
    \draw[->, PrimaryDeep!50] (\x) -- ++(0, -1.2);

% عنوان
\node[font=\normalsize\bfseries, PrimaryDeep] at (0, 4.5) {\rl{طیف مداخله: از فشار نرم تا زور نظامی}};

\end{tikzpicture}
\caption{طیف اشکال فشار/مداخله‌ی خارجی و سطح ریسک هر یک}
\label{fig:ch15-intervention-spectrum}
\end{figure}

\textbf{درس مهم:} بسیاری از بحث‌ها درباره‌ی «مداخله» فوراً به مداخله‌ی نظامی می‌پرند — در حالی که بیشتر طیف مداخله، غیرنظامی است. \concept{فشار خارجی غیرنظامی} (تحریم‌ها، حمایت از جامعه‌ی مدنی، فشار دیپلماتیک) هزینه‌ی بسیار کمتری دارد و — مانند مورد آپارتاید — می‌تواند بسیار مؤثر باشد. مشکل اصلی نجات‌طلبی فداکارانه آن است که معمولاً مستقیماً به انتهای طیف (مداخله‌ی نظامی) می‌پرد — بدون بررسی جدی گزینه‌های میانی.


%% ================================================================
\section{مسئله‌ی مشروعیت: چه کسی «حق» مداخله دارد؟}
\label{sec:ch15-legitimacy}
%% ================================================================

\needspace{6\baselineskip}
حتی اگر بپذیریم که نجات‌طلبی فداکارانه در موارد استثنایی موجه است، پرسش حقوقی-فلسفی مهمی باقی می‌ماند: \textbf{چه کسی حق مداخله دارد؟}

\subsection{دکترین «مسئولیت حمایت» و محدودیت‌هایش}
\label{subsec:ch15-r2p}

\concept{دکترین مسئولیت حمایت} (\lat{Responsibility to Protect} / \lat{R2P}) — مصوب اجلاس سران سازمان ملل در ۲۰۰۵ — سه رکن دارد:\footnote{\lr{United Nations General Assembly. "2005 World Summit Outcome," A/RES/60/1, 24 October 2005, paras. 138–140.}}

\begin{enumerate}
\item هر دولتی مسئولیت حمایت از شهروندانش در برابر نسل‌کشی، جنایات جنگی، پاک‌سازی قومی، و جنایات علیه بشریت را دارد.
\item جامعه‌ی بین‌المللی وظیفه دارد به دولت‌ها در ایفای این مسئولیت کمک کند.
\item اگر دولتی آشکارا ناتوان از حمایت از شهروندانش باشد یا خود عامل جنایت باشد، جامعه‌ی بین‌المللی حق مداخله دارد — از طریق شورای امنیت.

\end{enumerate}

\begin{tcolorbox}[enemybox, title={مشکلات عملی دکترین R2P}]
\begin{itemize}
\item \textbf{حق وتو:} هر یک از پنج عضو دائم شورای امنیت (آمریکا، روسیه، چین، فرانسه، بریتانیا) می‌تواند هر قطعنامه‌ای را وتو کند. روسیه و چین معمولاً مانع مداخله می‌شوند — حتی در مواردی مانند سوریه.
\item \textbf{اِعمال گزینشی:} R2P در لیبی اجرا شد اما در سوریه نه. در کوزوو اجرا شد اما در یمن نه. در رواندا اصلاً اجرا نشد. این گزینشی‌بودن، مشروعیت کل دکترین را زیر سؤال می‌برد.
\item \textbf{سوءاستفاده:} مداخله‌ی ناتو در لیبی (۲۰۱۱) با مجوز «حمایت از غیرنظامیان» آغاز شد اما به «تغییر رژیم» تبدیل شد — چیزی که مجوز شورای امنیت نداشت. این سوءاستفاده اعتماد روسیه و چین به R2P را از بین برد.\footnote{\lr{Bellamy, Alex J. "Libya and the Responsibility to Protect: The Exception and the Norm," \textit{Ethics \& International Affairs}, Vol. 25, No. 3, 2011, pp. 263–269.}}
\end{itemize}
\end{tcolorbox}

\subsection{جایگزین: «حق درخواست کمک» به جای «حق مداخله»}
\label{subsec:ch15-right-to-request}

شاید به جای «حق مداخله» (که ذاتاً از بالا به پایین و پدرسالارانه است)، باید از \concept{حق درخواست کمک} (\lat{Right to Request Assistance}) سخن گفت:

\begin{tcolorbox}[defbox, title={چارچوب پیشنهادی: حق درخواست کمک}]
\begin{enumerate}
\item مردم تحت ستم حق دارند از جامعه‌ی بین‌المللی \textit{درخواست} کمک کنند — اما این درخواست باید از کانال‌های مشخص و قابل‌تأیید باشد (نه فقط اپوزیسیون تبعیدی).
\item جامعه‌ی بین‌المللی \textit{وظیفه} دارد به درخواست پاسخ دهد — اما شکل پاسخ باید \concept{متناسب} و \concept{تدریجی} باشد (اول فشار دیپلماتیک، سپس تحریم، سپس حمایت مدنی، و فقط در آخرین مرحله مداخله‌ی نظامی).
\item مداخله‌ی نظامی فقط با مجوز شورای امنیت یا — در صورت انسداد شورای امنیت — با حمایت \concept{اکثریت بزرگ} (\lat{Supermajority}) مجمع عمومی مجاز است.
\item هر مداخله‌ای باید \textit{از پیش} شامل برنامه‌ی بازسازی، سازوکار خروج، و نظارت بین‌المللی باشد.
\end{enumerate}
\end{tcolorbox}


%% ================================================================
\section{آزمایش فکری: اگر شما تصمیم‌گیرنده بودید}
\label{sec:ch15-thought-experiment}
%% ================================================================

\needspace{6\baselineskip}
برای اینکه پیچیدگی این مسئله بهتر درک شود، خواننده را به یک آزمایش فکری دعوت می‌کنیم:

\begin{tcolorbox}[wavebox, title={آزمایش فکری: کشور X}]
فرض کنید:
\begin{itemize}
\item کشور X دیکتاتوری مطلق دارد. ۴۰ سال است هیچ تغییری رخ نداده.
\item رژیم سالانه حدود ۵٬۰۰۰ نفر را می‌کشد و ۵۰٬۰۰۰ نفر را زندانی/شکنجه می‌کند.
\item سرمایه‌ی انسانی در حال نابودی است: ۳۰٪ نخبگان فرار کرده‌اند.
\item ۶۰٪ جمعیت زیر خط فقر زندگی می‌کنند.
\item هیچ مقاومت مدنی سازمان‌یافته‌ای وجود ندارد (سرکوب کامل).
\item یک قدرت خارجی پیشنهاد مداخله می‌دهد. احتمال موفقیت (بر اساس بهترین تخمین): ۳۵٪.
\item هزینه‌ی مداخله: ۵۰٬۰۰۰–۱۰۰٬۰۰۰ کشته و ۵–۱۰ سال بی‌ثباتی.
\item هزینه‌ی عدم مداخله: تداوم وضع موجود به‌مدت نامعلوم.
\end{itemize}

\textbf{پرسش ۱:} آیا مداخله را می‌پذیرید؟

\textbf{پرسش ۲:} آیا پاسخ‌تان تغییر می‌کند اگر بدانید \textit{شما} جزو ۵۰٬۰۰۰ کشته‌ی احتمالی هستید؟

\textbf{پرسش ۳:} آیا پاسخ‌تان تغییر می‌کند اگر بدانید \textit{فرزندان‌تان} از آزادی پس از مداخله بهره‌مند خواهند شد؟

\textbf{پرسش ۴:} آیا پاسخ‌تان تغییر می‌کند اگر بدانید قدرت مداخله‌گر اهداف ژئوپولیتیک خودش را نیز دنبال می‌کند؟

\textbf{پرسش ۵:} آیا پاسخ‌تان تغییر می‌کند اگر احتمال موفقیت از ۳۵٪ به ۱۵٪ کاهش یابد؟
\end{tcolorbox}

این آزمایش فکری نشان می‌دهد که پاسخ به مسئله‌ی نجات‌طلبی فداکارانه عمیقاً \concept{وابسته به زمینه} (\lat{Context-Dependent}) است. هیچ پاسخ کلی‌ای — نه «همیشه آری» و نه «همیشه نه» — کافی نیست.


%% ================================================================
\section{موضع نهایی نویسنده: نه رد مطلق، نه پذیرش مطلق}
\label{sec:ch15-author-position}
%% ================================================================

\needspace{6\baselineskip}

\begin{tcolorbox}[mybox, title={موضع نهایی: واقع‌گرایی انتقادی}]
پس از بررسی قوی‌ترین استدلال‌ها به نفع و علیه نجات‌طلبی فداکارانه، نویسنده به موضعی می‌رسد که آن را \concept{واقع‌گرایی انتقادی} (\lat{Critical Realism}) می‌نامد:

\begin{enumerate}
\item \textbf{نجات‌طلبی فداکارانه به‌عنوان اصل عام مردود است.} شواهد تاریخی نشان می‌دهد که اکثر مداخلات خارجی شکست می‌خورند و وضع را بدتر می‌کنند. تبدیل نجات‌طلبی به «استراتژی» خطرناک است.

\item \textbf{اما رد مطلق مداخله نیز غیراخلاقی است.} رواندا، هولوکاست، و کامبوج نشان می‌دهند که عدم مداخله نیز می‌تواند جنایت باشد. موضع «هرگز مداخله نکنید» به‌اندازه‌ی «همیشه مداخله کنید» خطرناک است.

\item \textbf{هر مورد باید جداگانه ارزیابی شود} — بر اساس هفت شرط ذکرشده در بخش~\ref{sec:ch15-conditions}. هیچ قاعده‌ی کلی‌ای برای تمام موارد وجود ندارد.

\item \textbf{اولویت همواره با راه‌حل‌های غیرنظامی است:} فشار دیپلماتیک → تحریم‌ها → حمایت از جامعه‌ی مدنی → و \textit{فقط در آخرین مرحله} مداخله‌ی نظامی.

\item \textbf{عاملیت داخلی اصل است، حمایت خارجی فرع.} حتی در مواردی که کمک خارجی لازم است، باید \textit{ابزار} تقویت عاملیت داخلی باشد — نه جایگزین آن. مدل آفریقای جنوبی (عاملیت داخلی + فشار خارجی) بسیار موفق‌تر از مدل عراق (مداخله‌ی خارجی بدون عاملیت داخلی) بوده است.

\item \textbf{بدبینی سالم ضروری است:} هر ادعای «این بار متفاوت است» باید با شک روش‌مند بررسی شود. تجربه نشان داده که قدرت‌ها تقریباً همیشه منافع خود را بر منافع ملت مداخله‌شده ترجیح می‌دهند.

\item \textbf{مسئولیت نسلی واقعی است — اما تفسیرش پیچیده است.} نسل فعلی وظیفه‌ای نسبت به آیندگان دارد — اما این وظیفه ممکن است بهتر از طریق ساختن نهادهای مدنی (هرچند آهسته و ناقص) ایفا شود تا از طریق پذیرش مداخله‌ی خارجی پرخطر.
\end{enumerate}
\end{tcolorbox}


%% ================================================================
\section{پاسخ به پرسش آغازین فصل}
\label{sec:ch15-answering-question}
%% ================================================================

\needspace{6\baselineskip}
این فصل با پرسشی آغاز شد: \textbf{«وقتی بالاتر از سیاهی رنگی نیست، آیا هر تغییری بهتر از وضع موجود نیست؟»}

پاسخ — همان‌طور که تحلیل این فصل نشان داد — نه ساده و نه قطعی است:

\begin{tcolorbox}[quotebox, title={پاسخ سه‌لایه}]
\textbf{لایه‌ی اول — اخلاقی:} بله، وقتی انسان‌ها در نسل‌کشی، بردگی، یا سرکوب مطلق هستند، هر اقدامی برای نجات‌شان اخلاقاً قابل دفاع است. مسئولیت نسبت به بقای انسان‌ها بالاترین الزام اخلاقی است.

\textbf{لایه‌ی دوم — عملی:} اما «اخلاقاً قابل دفاع» به معنای «عملاً مؤثر» نیست. اکثر مداخلات خارجی نه تنها وضع را بهتر نکرده‌اند، بلکه بدتر کرده‌اند. نیت خوب + اجرای بد = فاجعه.

\textbf{لایه‌ی سوم — استراتژیک:} سؤال درست نه «آیا مداخله لازم است؟» بلکه «چه نوع مداخله‌ای، تحت چه شرایطی، توسط چه کسی، و با چه تعهدی؟» است. پاسخ به این سؤال در هر مورد متفاوت است — و هیچ فرمول کلی‌ای و
پاسخ به این سؤال در هر مورد متفاوت است — و هیچ فرمول کلی‌ای وجود ندارد.

\textbf{و مهم‌تر از همه:} «بالاتر از سیاهی رنگی نیست» فرضی نادرست است. تاریخ بارها نشان داده که بالاتر از سیاهی، سیاهی‌های عمیق‌تری وجود دارد — جنگ داخلی، تجزیه، ظهور گروه‌های تروریستی، فروپاشی کامل دولت، و نسل‌کشی‌های جدید. مداخله‌ی ناموفق می‌تواند تاریکی را نه روشن‌تر، بلکه بسیار تاریک‌تر کند.
\end{tcolorbox}


%% ================================================================
\section{درس‌های عملی: چه باید کرد؟}
\label{sec:ch15-practical-lessons}
%% ================================================================

\needspace{6\baselineskip}
اگر نه نجات‌طلبی مطلق و نه انزواگرایی مطلق پاسخ درست نیست، پس راه سوم چیست؟ این بخش تلاش می‌کند مجموعه‌ای از درس‌های عملی را — هم برای ملت‌های تحت سرکوب و هم برای جامعه‌ی بین‌المللی — استخراج کند.

\subsection{درس‌هایی برای ملت‌های تحت سرکوب}
\label{subsec:ch15-lessons-oppressed}

\begin{tcolorbox}[successbox, title={هفت درس عملی برای ملت‌های تحت سرکوب}]
\begin{enumerate}
\item \textbf{هرگز تمام تخم‌مرغ‌ها را در سبد مداخله‌ی خارجی نگذارید.} حتی اگر کمک خارجی بخشی از استراتژی باشد، بدون عاملیت داخلی قوی، هر مداخله‌ای محکوم به شکست است. مدل آفریقای جنوبی — عاملیت داخلی اول، فشار خارجی دوم — موفق‌ترین الگوی تاریخی است.

\item \textbf{هزینه‌ی واقعی مداخله را بشناسید — نه هزینه‌ی تبلیغاتی.} رسانه‌ها و گروه‌های لابی‌گر مداخله را «سریع، کم‌هزینه، و مؤثر» نشان می‌دهند. واقعیت تقریباً همیشه عکس این است. از عراقی‌ها، لیبیایی‌ها، و افغان‌ها بپرسید.

\item \textbf{نقشه‌ی «روز بعد» مهم‌تر از نقشه‌ی «روز سقوط» است.} سقوط دیکتاتور آسان‌ترین بخش کار است. ساختن نظم جدید — دهه‌ها طول می‌کشد. اگر نقشه‌ای برای روز بعد ندارید، سقوط رژیم ممکن است به فاجعه‌ای بدتر منجر شود.

\item \textbf{از انگیزه‌های مداخله‌گر آگاه باشید — و آن‌ها را مدیریت کنید.} هیچ قدرتی «رایگان» کمک نمی‌کند. قیمت مداخله چیست؟ پایگاه نظامی؟ قراردادهای نفتی؟ وابستگی سیاسی؟ امتیازات ژئوپولیتیک؟ بدانید چه می‌پردازید.

\item \textbf{اتحاد داخلی قبل از درخواست کمک خارجی.} اگر اپوزیسیون خود متفرق و متخاصم است، هیچ مداخله‌ی خارجی‌ای نمی‌تواند وحدت ایجاد کند. مداخله‌گر خارجی معمولاً شکاف‌های داخلی را تشدید می‌کند — نه التیام.

\item \textbf{ظرفیت‌سازی نهادی حتی در زیر سرکوب ممکن است.} لهستان دهه‌ی ۱۹۸۰ (جنبش همبستگی / \lat{Solidarność}) نشان داد که حتی در زیر سرکوب کمونیستی، می‌توان نهادهای موازی ساخت — سندیکاها، رسانه‌های زیرزمینی، شبکه‌های آموزشی.\footnote{\lr{Ost, David. \textit{Solidarity and the Politics of Anti-Politics}. Philadelphia: Temple University Press, 1990.}} این نهادها پس از سقوط رژیم، زیرساخت دموکراسی شدند.

\item \textbf{در بدترین حالت — نسل‌کشی یا فاجعه‌ی قریب‌الوقوع — درخواست کمک نه تنها حق بلکه وظیفه است.} اما این «بدترین حالت» واقعاً باید بدترین حالت باشد — نه بهانه‌ای برای نجات‌طلبی زودهنگام.
\end{enumerate}
\end{tcolorbox}

\subsection{درس‌هایی برای جامعه‌ی بین‌المللی}
\label{subsec:ch15-lessons-international}

\begin{tcolorbox}[wavebox, title={پنج درس برای جامعه‌ی بین‌المللی}]
\begin{enumerate}
\item \textbf{«مداخله نکردن» نیز یک تصمیم است — و باید پاسخگوی پیامدهایش باشید.} رواندا ۱۹۹۴ نشان داد که سکوت در برابر نسل‌کشی، همدستی است. جامعه‌ی بین‌المللی نمی‌تواند بگوید «به ما ربطی ندارد» و سپس ادعای ارزش‌های جهان‌شمول حقوق بشر کند.

\item \textbf{فشار زودهنگام، غیرنظامی، و پایدار مؤثرتر از مداخله‌ی نظامی دیرهنگام است.} تحریم‌های هدفمند، حمایت از جامعه‌ی مدنی، فشار دیپلماتیک، و انزوای بین‌المللی — اگر به‌موقع و با تداوم اعمال شوند — می‌توانند از رسیدن وضعیت به «انسداد مطلق» جلوگیری کنند. پیشگیری ارزان‌تر و مؤثرتر از درمان است.\footnote{\lr{Carnegie Commission on Preventing Deadly Conflict. \textit{Preventing Deadly Conflict: Final Report}. Washington, DC: Carnegie Corporation of New York, 1997.}}

\item \textbf{اگر مداخله‌ی نظامی ناگزیر شد، بدون تعهد بلندمدت بازسازی آغاز نکنید.} «مأموریت انجام شد» (\lat{Mission Accomplished}) — جمله‌ی مشهور جورج بوش در ۲۰۰۳ — نمونه‌ی بارز توهمی است که مداخله‌ی نظامی = پایان مشکل. مداخله بدون بازسازی نهادی، اقتصادی، و اجتماعی، فقط ویرانی تولید می‌کند.

\item \textbf{از مداخله‌ی یک‌جانبه بپرهیزید.} مداخله‌ی چندجانبه (با حمایت سازمان ملل، سازمان‌های منطقه‌ای، و افکار عمومی جهانی) بسیار بیشتر از مداخله‌ی یک‌جانبه (توسط یک قدرت) شانس موفقیت دارد — هم به‌دلیل تقسیم هزینه‌ها و هم به‌دلیل مشروعیت بیشتر.

\item \textbf{اصلاح ساختار شورای امنیت ضروری است.} تا زمانی که حق وتو وجود دارد، اجرای منسجم R2P غیرممکن است. پیشنهاد \concept{«مسئولیت وتو‌نکردن»} (\lat{Responsibility Not to Veto}) — مطرح‌شده توسط فرانسه و مکزیک — گامی در جهت درست است: اعضای دائم شورای امنیت باید متعهد شوند در موارد نسل‌کشی و جنایات علیه بشریت از حق وتو استفاده نکنند.\footnote{\lr{Luck, Edward C. "The Responsibility to Protect: The First Decade," \textit{Global Responsibility to Protect}, Vol. 3, No. 4, 2011, pp. 387–399.}}
\end{enumerate}
\end{tcolorbox}


%% ================================================================
\section{فراتر از دوگانه: به‌سوی «رهایی ترکیبی»}
\label{sec:ch15-hybrid-liberation}
%% ================================================================

\needspace{6\baselineskip}
دوگانه‌ی «عاملیت داخلی یا مداخله‌ی خارجی» یک دوگانه‌ی کاذب است. واقعیت — همان‌طور که تاریخ نشان می‌دهد — بسیار پیچیده‌تر است. موفق‌ترین رهایی‌ها ترکیبی بوده‌اند:

\begin{tcolorbox}[defbox, title={مفهوم: رهایی ترکیبی}]
\concept{رهایی ترکیبی} (\lat{Hybrid Liberation}) مدلی است که در آن:
\begin{enumerate}
\item \textbf{عاملیت داخلی محور اصلی است:} جنبش‌های مدنی، نهادهای موازی، رهبری بومی، و سازماندهی اجتماعی — ستون فقرات تغییر.
\item \textbf{فشار خارجی ابزار تکمیلی است:} تحریم‌ها، فشار دیپلماتیک، حمایت مالی و آموزشی از جامعه‌ی مدنی، رسانه‌های بین‌المللی — ابزارهایی که قدرت مقاومت داخلی را تقویت می‌کنند.
\item \textbf{مداخله‌ی نظامی آخرین راه‌حل است:} فقط در موارد نسل‌کشی یا فاجعه‌ی قریب‌الوقوع — و با تمام هفت شرط بخش~\ref{sec:ch15-conditions}.
\item \textbf{رابطه‌ی دوسویه، نه یک‌سویه:} ملت تحت سرکوب «مشتری» مداخله‌گر نیست — شریک اوست. تصمیمات اصلی باید توسط نهادهای بومی گرفته شود، نه توسط مداخله‌گر خارجی.
\end{enumerate}
\end{tcolorbox}

\textbf{نمونه‌های تاریخی رهایی ترکیبی:}

\begin{itemize}
\item \textbf{آفریقای جنوبی (۱۹۹۴):} مقاومت داخلی ANC + تحریم‌های بین‌المللی + فشار افکار عمومی جهانی. عاملیت داخلی محور بود؛ فشار خارجی تکمیلی.
\item \textbf{لهستان (۱۹۸۹):} جنبش همبستگی + حمایت پاپ ژان پل دوم + فشار غرب + تغییرات درون شوروی (پروسترویکای گورباچف). ترکیب عوامل داخلی و خارجی.
\item \textbf{تونس (۲۰۱۱):} انقلاب مردمی + حمایت (محدود) بین‌المللی + سنت نهادی پیشین (سندیکاها و جامعه‌ی مدنی نسبتاً قوی). تنها «موفقیت» بهار عربی.\footnote{\lr{Stepan, Alfred. "Tunisia's Transition and the Twin Tolerations," \textit{Journal of Democracy}, Vol. 23, No. 2, 2012, pp. 89–103.}}
\item \textbf{هند (۱۹۴۷):} مقاومت مدنی گاندی + فشار اقتصادی جنگ جهانی دوم بر بریتانیا + تغییر افکار عمومی بریتانیا + فشار آمریکا بر بریتانیا برای ترک مستعمرات.
\end{itemize}

\begin{tcolorbox}[wavebox, title={الگوی رهایی ترکیبی: چهار مرحله}]
\textbf{مرحله‌ی ۱ — سازماندهی داخلی:}
ساختن نهادهای مدنی موازی (سندیکاها، شبکه‌های آموزشی، رسانه‌های زیرزمینی، صندوق‌های کمک متقابل). حتی در سخت‌ترین شرایط سرکوب، سطحی از سازماندهی ممکن است.

\textbf{مرحله‌ی ۲ — بین‌المللی‌سازی:}
جلب توجه جهانی به وضعیت سرکوب از طریق رسانه‌ها، سازمان‌های حقوق بشری، دیاسپورا، و لابی‌گری دیپلماتیک. هدف: تبدیل مسئله از «مسئله‌ی داخلی» به «مسئله‌ی بین‌المللی».

\textbf{مرحله‌ی ۳ — فشار تدریجی:}
افزایش تدریجی فشار خارجی: اول بیانیه‌ها → سپس تحریم‌های هدفمند → سپس انزوای دیپلماتیک → سپس تحریم‌های جامع. هر مرحله فرصتی به رژیم می‌دهد تا اصلاح کند — و عدم اصلاح، فشار بیشتر را موجه می‌کند.

\textbf{مرحله‌ی ۴ — لحظه‌ی تغییر:}
وقتی ترکیب فشار داخلی و خارجی به نقطه‌ی بحرانی رسید، تغییر رخ می‌دهد — از طریق مذاکره (آفریقای جنوبی)، انقلاب مردمی (تونس)، یا فروپاشی درونی (شوروی). مداخله‌ی نظامی فقط اگر رژیم به نسل‌کشی یا سرکوب فاجعه‌بار متوسل شود.
\end{tcolorbox}


%% ================================================================
\section{نتیجه‌گیری: تراژدی انتخاب}
\label{sec:ch15-conclusion}
%% ================================================================

\needspace{6\baselineskip}
این فصل عمداً سخت‌ترین پرسش این کتاب را مطرح کرد — پرسشی که نویسنده نیز پاسخ قطعی‌ای برایش ندارد. این صداقت لازم است: هر کسی که ادعا کند پاسخ قطعی دارد — چه به نفع مداخله و چه علیه آن — یا مسئله را ساده‌سازی کرده، یا از تجربه‌ی واقعی سرکوب و جنگ بی‌خبر است.

\begin{tcolorbox}[quotebox, title={کلام پایانی فصل}]
آنچه این فصل نشان داد:

\begin{enumerate}
\item \textbf{نجات‌طلبی فداکارانه استدلال‌های قوی‌ای دارد} — به‌ویژه در مورد نسل‌کشی و انسداد مطلق. نادیده‌گرفتن این استدلال‌ها، بی‌انصافی نسبت به ستمدیدگان است.

\item \textbf{اما شواهد تاریخی بیشتر علیه آن‌اند تا له.} اکثر مداخلات خارجی شکست خورده‌اند — و شکست‌ها فاجعه‌بارتر از وضع قبل بوده‌اند.

\item \textbf{تفاوت اصلی میان موفقیت و شکست:} نه در «اصل مداخله» بلکه در شرایط زمینه‌ای (سنت نهادی، تعهد بلندمدت، همگنی جامعه) و نحوه‌ی اجرا (چندجانبه/یک‌جانبه، با بازسازی/بدون بازسازی).

\item \textbf{راه سوم — رهایی ترکیبی — موفق‌ترین مدل تاریخی است:} عاملیت داخلی قوی + فشار خارجی هدفمند + مداخله‌ی نظامی فقط به‌عنوان آخرین گزینه.

\item \textbf{هیچ فرمول کلی‌ای وجود ندارد.} هر وضعیت منحصربه‌فرد است و نیاز به تحلیل مستقل دارد. «قاعده‌ی سرانگشتی» (\lat{Rule of Thumb}) ممکن نیست — فقط \concept{حکم عملی} (\lat{Practical Judgment}) در هر مورد خاص.
\end{enumerate}

و شاید مهم‌ترین درس این فصل:

\textbf{بهترین زمان برای جلوگیری از رسیدن به دوراهی «نجات‌طلبی یا تسلیم»، پیش از رسیدن به آن نقطه است.} سرمایه‌گذاری در نهادهای مدنی، آموزش، رسانه‌های مستقل، و شبکه‌های اجتماعی مقاوم — \textit{پیش از} بسته‌شدن کامل فضا — بهترین بیمه علیه انسداد مطلق است. هر ملتی که این سرمایه‌گذاری را انجام دهد، کمتر به دوراهی فلج‌کننده‌ی «صبر ابدی یا مداخله‌ی خارجی» می‌رسد.

\textit{و اگر، با وجود همه‌ی تلاش‌ها، به آن نقطه رسید؟}

آنگاه — و فقط آنگاه — نجات‌طلبی فداکارانه نه «آرزواندیشی» بلکه شاید تنها گزینه‌ی باقی‌مانده است. اما حتی در آن لحظه، باید با \textbf{چشمان باز}، \textbf{محاسبه‌ی سرد}، و \textbf{بدون توهم} عمل کرد — نه با امیدی رمانتیک به «ناجی‌ای که همه‌چیز را درست می‌کند.»

زیرا تاریخ یک چیز را بی‌رحمانه ثابت کرده است:

\textbf{ناجیان بیرونی، حتی وقتی واقعاً نجات می‌دهند، صورتحسابشان را ارائه خواهند کرد.}
\end{tcolorbox}

\vspace{1em}
\begin{center}
\rule{0.5\linewidth}{0.4pt}
\end{center}
\vspace{1em}

\begin{tcolorbox}[mybox, title={خلاصه‌ی فصل ۱۵ در یک نگاه}]
\begin{itemize}
\item نجات‌طلبی فداکارانه = پذیرش آگاهانه‌ی هزینه‌ی مداخله‌ی خارجی برای رهایی نسل‌های آینده
\item پنج چارچوب فلسفی بررسی شد: فایده‌گرایی، وظیفه‌گرایی، عدالت بین‌نسلی، واقع‌گرایی تراژیک، فلسفه‌ی رهایی
\item شواهد تاریخی: آلمان نازی، آپارتاید، رواندا، کوزوو، کامبوج
\item ده استدلال قوی به نفع نجات‌طلبی فداکارانه ارائه و بررسی شد
\item شش نقد اساسی: مغالطه‌ی تعمیم، مسئله‌ی نمایندگی، انگیزه‌ی مداخله‌گر، امکان بدترشدن وضع، مشکلات محاسبه‌ی فایده‌گرایانه، خطر ابزاری‌شدن
\item هفت شرط لازم برای مداخله‌ی موجه تعریف شد
\item مدل «رهایی ترکیبی» به‌عنوان جایگزین پیشنهاد شد: عاملیت داخلی + فشار خارجی + مداخله‌ی نظامی فقط در آخرین مرحله
\item موضع نهایی: واقع‌گرایی انتقادی — نه رد مطلق، نه پذیرش مطلق
\end{itemize}
\end{tcolorbox}